\documentclass[11pt]{article}
\usepackage{graphicx} % Required for inserting images
\setlength{\parindent}{0pt}
\usepackage[bookmarks=true, colorlinks=true, urlcolor=blue, linkcolor=black]{hyperref}

\usepackage{enumitem}
\usepackage[utf8]{inputenc}
\usepackage[T1]{fontenc}
\usepackage[english]{babel}
\usepackage{geometry}
\usepackage{lastpage}
\usepackage{fancyhdr}
\usepackage{datetime}
\usepackage{parskip}

% Link to other SOPs, forms
\usepackage{xr}
\externaldocument[301-]{../300_building/1_access/sop_301_access}
% Definitions for SOP sections 3.
\newcommand{\AdverseEvent}{\textbf{Adverse Event}: Any untoward or unfavorable medical occurrence in a human subject, including any abnormal sign, symptom, or disease, temporally associated with the subject's participation in the research, whether or not considered related to the subject's participation in the research.}

\newcommand{\ControlRoom}{\textbf{Control Room}: B17; The main room of the MRI suite that contains the computer displays and keyboards that control the MRI scanner. The Control Room is accessed from B15 that is controlled by an Electronic Lock.}

\newcommand{\DcmComputer}{\textbf{DICOM Computer}: The Mac Studio located in the Control Room the receives DICOM files from the scanner.}

\newcommand{\ElectronicLock}{\textbf{Electronic Lock}: All doors at the ND-HNC are controlled by electronic locks. The electronic lock system recognizes the owner's Irish1Card, logs the time of the attempted entry, and unlocks the door if the owner has been granted access.}

\newcommand{\EEG}{\textbf{EEG}: Electroencephalograpy, a technique used to measure changes in electromagnetic properties of neural tissue.}

\newcommand{\EquipmentRoom}{\textbf{Equipment Room}: B18; The room containing hardware for MRI operations.}

\newcommand{\FDA}{\textbf{FDA}: Food and Drug Administration.}

\newcommand{\GoodCatch}{\textbf{Good Catch}: Early recognition of an event with the potential to cause damage, injury, or illness.}

\newcommand{\Helper}{\textbf{Helper}: Second person in Zones 3 and 4 who assists in emergency responses, typically Level 2-3.}

\newcommand{\Incident}{\textbf{Incident}: An occurence, action, or event that warrants reporting (e.g. a policy is not followed and/or an individual's health is at risk). Includes Adverse Events, Near Misses, Unanticipated Problems, etc.}

\newcommand{\Lead}{\textbf{Lead}: Primary person in Zones 3 and 4 responsible for the Visitor's safety, in charge during emergency responses. Level 4 operator or MRI Technologist.}

\newcommand{\MagnetRoom}{\textbf{Magnet Room}: B17A; The room within the MRI suite that contains the MRI scanner. It is accessed from within the MRI Control room (B17) through a doorway that is controlled by an Electronic Lock.}

\newcommand{\NearMiss}{\textbf{Near Miss}: An incident that occurred but did not cause personal injury/illness e.g. a small/contained fire, trip on sidewalk or carpeting, falling object.}

\newcommand{\NDHNC}{\textbf{ND-HNC}: Notre Dame Human Neuroimaging Center, Lower Level of Veldman Family Psychology Clinic.}

\newcommand{\Operator}{\textbf{Operator}: MRI Technologist or Level 4 researcher, responsible for conducting research operations and maintaining equpiment. The Lead during emergency responses.}

\newcommand{\PIN}{\textbf{PIN}: Personal Identification Number, separate for Irish1Card and security.}

\newcommand{\ResearchAssistant}{\textbf{Research Assistant}: Individual with access Levels 1-3 to the MRI research suite. Typically the Helper during emergency responses.}

\newcommand{\SeriousInjury}{\textbf{Serious Injury}: An injury or illness that is life threatening; results in permanent impairment of a body function or permanent damage to a body structure; or necessitates medical or surgical intervention to preclude permanent damage or impairment.}

\newcommand{\StimControl}{\textbf{Stimulus Control Computer}: The Windows machines in the Control Room that deploy fMRI tasks and capture participant responses.}

% \newcommand{\Subject}{\textbf{Subject}: An individual who is participating in an experimental protocol at ND-HNC and has not taken the ND-HNC safety course.}

\newcommand{\Supervision}{\textbf{Supervision}: A supervisor is continuously present in both sight and sound range of the individual or task being monitored. This approach ensures the supervisor can actively observe, direct, and immediately intervene if a complication or safety issue arises.}

\newcommand{\TMS}{\textbf{TMS}: Transcranial magnetic stimulation, a technique used to disrupt normal operations of neural tissue.}

\newcommand{\TornadoWarn}{\textbf{Tornado Warning}: A tornado has been sighted or indicated by weather radar and there is imminent danger to life and property. Action is required.}

\newcommand{\TornadoWatch}{\textbf{Tornado Watch}: Tornados are possible in and near the watch area. Be ready to act quickly and move to the designated tornado safe area.}

\newcommand{\UnanticipatedProblem}{\textbf{Unanticipated Problem}: An incident, experience, or outcome that meets \textbf{all} of the following criteria: (1) unexpected (in terms of nature, severity, or frequency) given (a) the research procedures that are described in the protocol-related documents, such as the IRB-approved research protocol and informed consent document; and (b) the characteristics of the subject population being studied; (2) related or possibly related to participation in the research; and (3) suggests that the research places subjects or others at a greater risk of harm than was previously known or recognized.}

\newcommand{\VFPC}{\textbf{VFPC}: Veldman Family Psychology Clinic.}

\newcommand{\Visitor}{\textbf{Visitor}: An individual who has not taken the ND-HNC MRI Safety Course but has a legitimate reason for entering the MRI suites. Research participants and their legal guardians are Visitors. Other Visitors may include: individuals providing/receiving a tour of the VFPC and ND-HNC, faculty and staff who have not received MRI Safety Course, and research assistants.}



% figures and tables
\usepackage{graphicx}
\graphicspath{ {../../pics} }
\usepackage{rotating}
\usepackage{float}
\usepackage{subfig}
\usepackage{caption}
\DeclareCaptionLabelFormat{adja-page}{\hrulefill\\#1 #2 \emph{(previous page)}}


\renewcommand{\headrulewidth}{0pt}
% \newdateformat{yearmonth}{\THEYEAR-\THEMONTH-\THEDAY}
\pagestyle{fancy}
\fancyhead[L]{}
\fancyhead[R]{Approved: TBD}
\fancyfoot[C]{}
\fancyfoot[L]{ND-HNC SOP 202: v.0.1.0}
\fancyfoot[R]{p. \thepage /\pageref*{LastPage}}



% start document
\begin{document}

\renewcommand{\labelenumii}{\arabic{enumi}.\arabic{enumii}}
\renewcommand{\labelenumiii}{\arabic{enumi}.\arabic{enumii}.\arabic{enumiii}}
\renewcommand{\labelenumiv}{\arabic{enumi}.\arabic{enumii}.\arabic{enumiii}.\arabic{enumiv}}

\begin{center}
	\textbf{\Large SOP 202: Zone Access}\\
	\hrulefill
\end{center}


\section{Summary}
This document establishes MRI Safety Zones of the research suite at the Notre Dame Human Neuroimaging Center (ND-HNC). These Zones follow the American College of Radiology guidelines and are designed to ensure that Zone access is commensurate with training and experience. Roles and restrictions for each level of access are also described.


\section{Scope}
These policies apply to all individuals who desire entry into the research MRI suite (Fig. \ref{fig:mri-zones}).

\section{Definitions}

\begin{itemize}
	\item \ControlRoom
	\item \ElectronicLock
	\item \MagnetRoom
	\item \Subject
	\item \Visitor
\end{itemize}

\section{Notes}

\begin{enumerate}
	\item The MRI simulator (mock) is \textbf{NOT} compatible with pacemakers!
\end{enumerate}

\section{Policies \& Procedures}

In the sections below Safety Zones, user access, and user responsibilities are described.


% research suite
\subsection{Research Suite}
\label{ssec:res-suite}

The ND-HNC research suite consists of the following rooms and areas:

\begin{enumerate}
	\item Reception
	\item Staff offices
	\item Data processing and conference
	\item Interview and changing
	\item EEG: preparation, testing, and control
	\item TMS: administration
	\item MRI: simulator
	\item MRI: vestibule, equipment, control
	\item MRI: magnet room
\end{enumerate}

These are located on the lower level of the Veldman Family Psychology Clinic. See SOP 301 for access procedures.

[TODO PIC]


% Zones
\subsection{Safety Zones}
\label{ssec:safety-zones}

The research suite is organized into four Safety Zones. [TODO PIC]

\subsubsection{Zone 1}
\label{sssec:safety-zones-1}

This area constitutes the reception area, staff offices, conference room, and processing room (\ref{ssec:res-suite}, 1-3). The reception area is accessible from the building's main reception and is used as a waiting area for researchers and participants. A bathroom and two interview rooms are accessible from the waiting area. The staff offices, conference room, and processing room are accessible from the staff entrance and are used for ND-HNC related administration, research, and meetings. A break room and staff bathroom are also located in this zone.

MRI compatibility screening is \textbf{NOT} required for entry into Zone 1, and is typically the location where such screening is conducted. [TODO PIC]

\subsubsection{Zone 2}
\label{sssec:safety-zones-2}

This area consists of changing rooms, bathroom and shower, interview rooms, MRI simulator, EEG, and TMS rooms (\ref{ssec:res-suite}, 4-7). Zone 2 is accessed from Zone 1 with the Irish1Card. Zone 2 is used for research protocols that do not involve the MRI machine (e.g. behavioral testing, clinical assessment, EEG, and TMS). MRI compatibility screening can be conducted in Zone 2, but is typically done in Zone 1. Participants and visitors must be escorted in this Zone by a researcher with Level 1 access.

MRI compatibility screening is \textbf{NOT} required for entry into Zone 2 when the participant will not contribute MRI data.


\subsubsection{Zone 3}
\label{sssec:safety-zones-3}

This area consists of the MRI vestibule, equipment, and control rooms (\ref{ssec:res-suite}, 8). Zone 3 is accessed from Zone 2 with the Irish1Card. This area is used for collecting MRI data and maintaining MR-related hardware. Researchers with Level 2 access can enter this area for MRI-related purposes, and participants and visitors must be escorted by a researcher with Level 2 access.

MRI compatibility screening \textbf{IS} required for entry into Zone 3, and the screening form must have been reviewed and approved by a researcher with Level 2 access.


\subsubsection{Zone 4}
\label{sssec:safety-zones-4}

This area consists of the Magnet Room (\ref{ssec:res-suite}, 9), where the MRI machine is located. Zone 4 is accessed from the Control Room (Zone 3). Researchers with Level 3 and 4 access must be supervised when entering this area by the MRI Technologist or a researcher with Level 5 access.

MRI compatibility screening \textbf{IS} required for entrance into Zone 3, and this screening must have been reviewed and approved by the MRI Technologist or a researcher with Level 4 access (in addition to the Level 2 approval noted above).


% Levels
\subsection{Access Levels}
\label{ssec:access-levels}

Access to the ND-HNC research suite is restricted to research personnel, participants, and legal guardians. Anyone with building access is able to visit Zone 1. Access to additional Zones requires a demonstrated and current need as well as additional training and certifications.

\subsubsection{Level 1}
\label{sssec:access-levels-1}

Level 1 researchers may enter Zone 2 in order to access the interview, EEG, and TMS rooms. They escort visitors and participants in these areas. Their level of access relates to non-MRI activities.

Requirements:
\begin{itemize}
	\item Complete annual ND-HNC Safety Training 1
	\item Complete annual first-aid training
	\item Participate in Maintenance Training 1
	\item Wear Level 1 access badge while in the ND-HNC research suite
\end{itemize}

Privileges:
\begin{itemize}
	\item Unsupervised access to Zone 2
	\item Escort participants and visitors in Zone 2
\end{itemize}

Restrictions:
\begin{itemize}
	\item No access to Mock Room (B14)
	\item No access to Zone 3 (B15, B17, B18)
	\item No access to Zone 4 (B17a)
\end{itemize}


\subsubsection{Level 2}
\label{sssec:access-levels-2}

Level 2 researchers may enter Zone 3 to assist in MRI-related research activities. They can access the Mock, Stimulus Control, and DICOM computers, and can review MRI compatibility screeners. Most research assistants in the MRI suite are Level 2. Their level of access relates to assisting MRI research.

Requirements:
\begin{itemize}
	\item Complete annual ND-HNC Safety Training 2
	\item Complete annual first-aid and CPR training
	\item Annually provide an MRI compatibility screener
	\item Participate in Maintenance Training 2
	\item Wear Level 2 access badge while in the ND-HNC research suite
\end{itemize}

Privileges:
\begin{itemize}
	\item All Level 1 privileges
	\item Access and use of Mock Scanner and Mock computer
	\item Supervised access to Zone 3
	\item Access Stimulus Control, DICOM computers
	\item Review, sign MRI compatibility screeners
\end{itemize}

Restrictions:
\begin{itemize}
	\item No access to Zone 3 Equipment Room
	\item No access to Zone 4
\end{itemize}


\subsubsection{Level 3}
\label{sssec:access-levels-3}

Level 3 researchers are more independent in Zone 3 and may access Zone 4 when invited. They are graduate students and post-doctoral fellows.

Requirements:
\begin{itemize}
	\item Complete annual ND-HNC Safety Training 3
	\item Complete annual first-aid and CPR training
	\item Participate in Maintenance Training 3
	\item Wear Level 3 access badge while in the ND-HNC research suite
	\item Current Notre Dame graduate student or post-doctoral fellow
	\item Complete 20 hours of research assistance in Zone 3 as Level 2
\end{itemize}

Privileges:
\begin{itemize}
	\item All Level 2 privileges
	\item Unsupervised access to Zone 3
	\item Supervised access to Zone 4
	\item Assist in preparing participant
\end{itemize}

Restrictions:
\begin{itemize}
	\item No access to Zone 3 Equipment Room
	\item No unsupervised access to Zone 4
	\item No access to MRI Control computer
\end{itemize}


\subsubsection{Level 4}
\label{sssec:access-levels-4}

Level 4 researchers are more independent in the MRI research suite, and may train to operate the MRI scanner for project-specific needs (e.g. scanning adolescents on the weekend, scanning adults at night). They are graduate students and post-doctoral fellows. Their access pertains to learning to operate, maintain, and troubleshoot the scanner.

Requirements:
\begin{itemize}
	\item Complete annual ND-HNC Safety Training 4
	\item Complete annual first-aid and CPR training
	\item Participate in Maintenance Training 4
	\item Wear Level 4 access badge while in the ND-HNC research suite
	\item Current Notre Dame graduate student or post-doctoral fellow
	\item Complete 30 hours of research assistance in Zones 3, 4 as Level 3
\end{itemize}

Privileges:
\begin{itemize}
	\item All Level 3 privileges
	\item Supervised access to MRI Control computer
	\item Supervised access to Zone 3 Equipment Room
	\item Perform QA/QC on phantoms
\end{itemize}

Restrictions:
\begin{itemize}
	\item Cannot be Lead in Emergency Response SOPs [TODO LINK]
\end{itemize}


\subsubsection{Level 5}
\label{sssec:access-levels-5}

Level 5 researchers are independent in the MRI research suite, but operate under guidance and with permission from the MRI Technologist.They can serve as Lead during scanning, and when working in that capacity they are the final check of participant MRI compatibility (and most not be the same as the first check). They are Ph.D. candidates, post-doctoral fellows, and faculty with a research appointment.

Requirements:
\begin{itemize}
	\item Pass operator exam
	\item Pass emergency SOP exam
	\item Participate in maintenance training
	\item Wear Level 5 access badge while in the ND-HNC research suite
	\item Current Notre Dame Ph.D. candidate, post-doctoral fellow, or faculty
	\item Complete 30 hours of training as Level 4
\end{itemize}

Privileges:
\begin{itemize}
	\item All Level 4 privileges
	\item Unsupervised access to Zones 3 and 4
	\item Screen and scan participants
	\item Perform start-up procedures
	\item Perform maintenance scans and troubleshoot issues
	\item Perform close-down procedures
	\item Maintain relevant logs
	\item Coordinate with ND-HNC staff
	\item Train Levels 2-4
\end{itemize}

Restrictions:
\begin{itemize}
	\item Must be accompanied by Level 2+ when a participant or visitor is in Zone 4
	\item Operates at the discretion of the MRI Technologist and other ND-HNC staff
\end{itemize}


\section{Figures}

% MRI zones
\begin{figure}[H]
	\centering
	\includegraphics[width=\textwidth]{mri_zones.jpg}
	\caption{Schematic of ND-HNC MRI Zones. \textbf{Orange} Zone 1. \textbf{Purple} Zone 2. \textbf{Green} Zone 3. \textbf{Blue} Zone 4.}
	\label{fig:mri-zones}
\end{figure}


\end{document}